\chapter{Marco teórico y estado del arte}

\section{Antecedentes y estado del arte}

La revisión que sigue es preliminar: reúne fundamentos de análisis espacial y de arquitectura de software para justificar el diseño y evaluación del artefacto propuesto. Deberá actualizarse con una búsqueda sistemática antes de la sustentación.

\subsection{Antecedentes internacionales}

Los trabajos recientes muestran que la concentración de incidentes en unos pocos lugares no es un fenómeno aislado ni exclusivo de contextos específicos. \citet{ristea2020crimegis} aplicaron sistemas de información geográfica al análisis urbano del delito; \citet{hajela2020hotspot} estudiaron la identificación de zonas críticas mediante agrupamiento; y \citet{masron2021crimehotspot} compararon la distribución de zonas críticas en dos distritos de Malasia. Para datos espaciotemporales, \citet{choi2021mdstdbscan} propusieron una variante de DBSCAN para clustering multidimensional, mientras que \citet{xie2021significantdbscan} incorporaron inferencia estadística al agrupamiento basado en densidad. Otros trabajos exploran clustering jerárquico y modelos difusos para cartografía y propagación de concentraciones delictivas \citep{gustriansyah2022crime,cardone2022fuzzyhotspot}. Estos antecedentes muestran que la selección del algoritmo debe responder al tipo de dato y que el resultado algorítmico debe distinguirse de la evidencia estadística.

Como referencias metodológicas recientes, \citet{medinaariza2023points} presentan métodos de patrones puntuales para eventos delictivos y \citet{medinaariza2023time} discuten la dimensión temporal del mapeo del delito; \citet{medinaariza2023hotspots} abordan la detección de concentraciones y repeticiones. Estos aportes orientan la separación entre descripción cartográfica, análisis de patrones y evaluación de concentraciones.

\subsection{Antecedentes nacionales}

\citet{cutid2025} analizaron el delito en la Provincia Constitucional del Callao (febrero de 2024) con minería de datos espaciales: emplearon densidad kernel para visualizar \textit{hotspots} y DBSCAN, con $\varepsilon$ seleccionado mediante el gráfico de $k$-distancias, para detectar grupos discretos. Reportan una fuerte concentración geográfica de los incidentes y un grupo de 317 ocurrencias en el Cercado del Callao, y concluyen que el delito sigue patrones geográficos definidos. Es el antecedente más cercano a este plan por la técnica y por el contexto peruano.

Un sistema de geo-sectorización de zonas delictivas con K-Means y visualización en mapas, aplicado a un contexto turístico, exige fijar el número de grupos de antemano y asigna todos los puntos a un grupo, a diferencia de DBSCAN, que admite ruido y formas arbitrarias.

\subsection{Antecedentes locales}

En la búsqueda preliminar no se identificaron estudios publicados que apliquen técnicas de agrupamiento espacial a registros de serenazgo del distrito de San Sebastián o de la ciudad del Cusco. Esta afirmación es provisional y se confirmará con una búsqueda en los repositorios institucionales (\gls{unsaac}, ALICIA y RENATI).

\subsection{Síntesis y brecha de investigación}

La Tabla~\ref{tab:estadoarte} resume los antecedentes sobre concentración y análisis espacial. Estos métodos aportan algoritmos para detectar grupos, pero no especifican por sí mismos cómo integrar las fuentes, validaciones, transformaciones, parámetros, pruebas y productos en un flujo reproducible. La brecha que aborda este plan es la falta de una arquitectura analítica documentada y evaluada para caracterizar los incidentes registrados en SIPCOP-M de San Sebastián, considerando sus restricciones de calidad, resolución espacial y cobertura temporal. La existencia de grupos o concentraciones significativas no se presupone.

\begin{table}[H]
\centering
\caption{Síntesis de los antecedentes revisados.}
\label{tab:estadoarte}
\begingroup\singlespacing\footnotesize
\begin{tabular}{@{}p{3.3cm}p{3.2cm}p{4.3cm}p{3.3cm}@{}}
\toprule
\textbf{Estudio} & \textbf{Técnica} & \textbf{Aporte} & \textbf{Limitación respecto de este plan} \\
\midrule
\citet{ristea2020crimegis} & Sistemas de información geográfica & Análisis urbano del delito y cartografía espacial & Enfoque general, no arquitectura analítica completa \\
\citet{hajela2020hotspot} & Agrupamiento para identificación de zonas críticas & Identificación de zonas críticas mediante clustering & Sin flujo de validación y trazabilidad \\
\citet{cutid2025} & Densidad kernel y DBSCAN & Aplicación en el Callao, $\varepsilon$ por $k$-distancias & Otra jurisdicción; un solo mes \\
\citet{gustriansyah2022crime} & Clustering jerárquico & Mapeo de tasas delictivas & Sin evaluación de reproducibilidad ni trazabilidad \\
\citet{cardone2022fuzzyhotspot} & Modelos difusos & Propagación de intensidad de hotspots & No especifica arquitectura de software \\
\bottomrule
\end{tabular}
\endgroup
\fuente{Elaboración propia.}
\end{table}

\section{Fundamentos teóricos}

\subsection{Seguridad ciudadana y concentración del delito}

La seguridad ciudadana designa la protección de las personas frente a hechos que amenazan su integridad y sus bienes, y en el Perú se organiza como un sistema con participación de los gobiernos locales. El \termref{term-serenazgo}{serenazgo} es el servicio municipal que realiza patrullaje y prevención en el distrito. La teoría de las actividades rutinarias sostiene que el delito se produce cuando coinciden en el espacio y el tiempo un ofensor motivado, un objetivo adecuado y la ausencia de guardianes; de ello se sigue que los incidentes tienden a agruparse en ciertos lugares y horarios, y que esa regularidad pueda estudiarse mediante mapas y estadística \citep{ristea2020crimegis,medinaariza2023hotspots}.

\subsection{Análisis espacial}

El análisis espacial estudia la distribución geográfica de los eventos, su densidad y su \termref{term-spatial-clustering}{agrupamiento} \citep{ristea2020crimegis}. Una vez confirmada la \termref{term-georeferencing}{georreferenciación} y que las coordenadas del archivo pueden transformarse al sistema \gls{wgs84}, la distancia entre dos ocurrencias se calculará con la fórmula \termref{term-haversine}{haversine}:
\begin{equation}
\symref{sym-distance}{d(p,q)} = 2\symref{sym-earth-radius}{R} \arcsin\sqrt{\sin^{2}\!\left(\frac{\symref{sym-latitude}{\varphi}_q-\symref{sym-latitude}{\varphi}_p}{2}\right) + \cos\symref{sym-latitude}{\varphi}_p \cos\symref{sym-latitude}{\varphi}_q \sin^{2}\!\left(\frac{\symref{sym-longitude}{\lambda}_q-\symref{sym-longitude}{\lambda}_p}{2}\right)}
\label{eq:haversine}
\eqtitulo{Distancia haversine entre dos ocurrencias.}
\end{equation}
donde $\symref{sym-earth-radius}{R}$ es el radio medio de la Tierra, $\symref{sym-latitude}{\varphi}$ la latitud y $\symref{sym-longitude}{\lambda}$ la longitud, ambas en radianes.

La densidad de eventos en una ubicación $x$ puede estimarse mediante \termref{term-kernel-density}{densidad kernel}:
\begin{equation}
\symref{sym-density}{\hat{f}(x)} = \frac{1}{n \symref{sym-bandwidth}{h}^{2}} \sum_{i=1}^{n} \symref{sym-kernel}{K(\cdot)}\!\left(\frac{x - x_i}{\symref{sym-bandwidth}{h}}\right)
\label{eq:kde}
\eqtitulo{Estimador de densidad kernel.}
\end{equation}
donde $n$ es el número de ocurrencias, $\symref{sym-bandwidth}{h}$ el ancho de banda y $\symref{sym-kernel}{K(\cdot)}$ la función kernel. La densidad kernel produce una superficie continua apta para \termref{term-heatmap}{mapas de calor}, pero no delimita grupos de ocurrencias; en este estudio se usa solo como visualización.

\subsection{Agrupamiento espacial por densidad}

\gls{dbscan} \citep{choi2021mdstdbscan,xie2021significantdbscan} agrupa los puntos según su densidad a partir de la vecindad de cada uno:
\begin{equation}
\symref{sym-neighborhood}{N_{\varepsilon}(p)} = \{\, q \in D \mid \symref{sym-distance}{d(p,q)} \leq \symref{sym-epsilon}{\varepsilon} \,\}
\label{eq:vecindad}
\eqtitulo{Vecindad $\varepsilon$ de un punto en DBSCAN.}
\end{equation}
Un punto es \emph{núcleo} si $|\symref{sym-neighborhood}{N_{\varepsilon}(p)}| \geq \symref{sym-minpts}{\mathrm{MinPts}}$; los puntos alcanzables desde un núcleo forman un \termref{term-spatial-cluster}{grupo espacial}, y los restantes se etiquetan como \termref{term-noise}{ruido}. El método no requiere fijar el número de grupos, admite formas arbitrarias y es sensible a $\symref{sym-epsilon}{\varepsilon}$ y MinPts. Una guía habitual para elegir $\varepsilon$ es el gráfico de $k$-distancias, donde se busca el ``codo'' de la curva ordenada de distancias al $k$-ésimo vecino. La variante jerárquica de agrupamiento por densidad ofrece una alternativa útil cuando la estructura de densidad no es uniforme y no se desea fijar el radio explícitamente.

\subsection{Resolución y error de georreferenciación}

Las coordenadas de los registros policiales suelen proceder de la geocodificación de direcciones o de puntos de referencia, y su error posicional afecta la detección de concentraciones; por eso se recomienda informar la proporción de registros georreferenciados y la precisión lograda. Cuando las coordenadas están redondeadas a una malla, varios registros comparten la misma ubicación y las distancias menores que el paso de la malla no son observables. En consecuencia, el radio $\varepsilon$ de DBSCAN no debe ser menor que ese paso, los puntos coincidentes se cuentan como registros distintos al evaluar MinPts y la simulación del modelo nulo debe reproducir la misma cuantización.

\subsection{Inferencia sobre concentraciones espaciales y espacio-temporales}

Un grupo devuelto por un algoritmo no demuestra por sí solo una concentración mayor que la esperable por azar. Para contrastarla se dispone de simulaciones Monte Carlo del tamaño de los grupos bajo un modelo nulo, y de enfoques complementarios de comparación espacial y temporal que permiten evaluar si la concentración excede lo esperado en un marco de referencia. En este plan el contraste principal es la simulación Monte Carlo descrita en la subsección ``Concentración espacial frente al azar'' del capítulo de metodología; la comparación con modelos nulos y la evolución temporal se consideran evidencias complementarias.

\subsection{Evaluación de la calidad del agrupamiento}

El coeficiente de silueta compara, para cada punto $i$, la distancia media a los puntos de su grupo, $a(i)$, con la distancia media al grupo vecino más cercano, $b(i)$:
\begin{equation}
\symref{sym-silhouette}{s(i)} = \frac{\symref{sym-between}{b(i)} - \symref{sym-within}{a(i)}}{\max\{\symref{sym-within}{a(i)},\, \symref{sym-between}{b(i)}\}}, \qquad -1 \leq \symref{sym-silhouette}{s(i)} \leq 1
\label{eq:silueta}
\eqtitulo{Coeficiente de silueta.}
\end{equation}
Valores cercanos a 1 indican grupos compactos y bien separados. Como DBSCAN deja puntos sin grupo, la silueta se calcula sobre los puntos asignados y se informa junto con el porcentaje de ruido y el número de grupos.

\subsection{Arquitectura analítica y evaluación del artefacto}

Una arquitectura de software describe las estructuras relevantes de un sistema, sus elementos, relaciones y principios que orientan su evolución; no equivale únicamente a un diagrama ni a la implementación de una aplicación. La norma ISO/IEC/IEEE 42010 establece conceptos y requisitos para describir arquitecturas mediante interesados, preocupaciones, puntos de vista y vistas \citep{iso42010_2022}. Para este estudio, la arquitectura será lógica y analítica: especificará módulos, responsabilidades, interfaces, flujo de datos y decisiones de procesamiento, sin comprometer un despliegue institucional.

La investigación en ciencia del diseño exige identificar explícitamente el artefacto y distinguirlo del problema y del método empleados para construirlo. \citet{weigand2021artifact} proponen una ontología para describir artefactos de investigación; \citet{weigand2023artifact} ofrecen criterios para identificar el artefacto en estudios de sistemas de información. Además, \citet{barata2023limitations} señalan la importancia de informar limitaciones en este tipo de investigación, y \citet{schaeffer2025evaluation} proponen seleccionar métodos de evaluación adecuados a los objetivos y al artefacto. Estos criterios respaldan que este plan describa la arquitectura, la prueba de concepto y sus límites como productos diferenciados.

La calidad del artefacto se evaluará respecto de atributos observables y criterios definidos antes de las pruebas. ISO/IEC 25010 proporciona un modelo de calidad de producto que organiza características y subcaracterísticas para especificar y evaluar sistemas y productos de software \citep{iso25010_2023}; \citet{argotti2024quality} ejemplifican cómo operacionalizar modelos de calidad ISO 25000 para evaluar software. En esta investigación se seleccionarán los atributos pertinentes al alcance: adecuación funcional, confiabilidad de ejecución, mantenibilidad modular y reproducibilidad de resultados. La evaluación combinará trazabilidad entre requisitos y componentes, casos funcionales con entradas válidas e inválidas, repetición controlada del flujo y comprobación de los productos generados.

El proceso de construcción seguirá una lógica de investigación en ciencia del diseño: identificar el problema, definir objetivos de solución, diseñar y desarrollar el artefacto, demostrar su uso y evaluarlo frente a los objetivos establecidos. El extracto SIPCOP-M servirá como caso demostrativo, mientras que los resultados de agrupamiento serán evidencia del funcionamiento de módulos analíticos y no una prueba autónoma de calidad arquitectónica.

\subsection{Análisis temporal y pruebas estadísticas}

Para contrastar si la frecuencia de ocurrencias por categoría temporal (turno, día de la semana, mes) difiere de la esperada se utiliza la prueba chi-cuadrado de Pearson. En la prueba de homogeneidad, $\symref{sym-null}{H_0}$ establece que las distribuciones comparadas coinciden y $\symref{sym-alternative}{H_1}$ que difieren:
\begin{equation}
\symref{sym-chi-square}{\chi^{2}} = \sum_{i=1}^{k} \frac{(\symref{sym-observed}{O_i} - \symref{sym-expected}{E_i})^{2}}{\symref{sym-expected}{E_i}}
\label{eq:chi2}
\eqtitulo{Estadístico chi-cuadrado de Pearson.}
\end{equation}
donde $\symref{sym-observed}{O_i}$ y $\symref{sym-expected}{E_i}$ son las frecuencias observada y esperada de la categoría $i$. La prueba es fiable cuando las frecuencias esperadas no son demasiado pequeñas. Su magnitud depende del tamaño de muestra, por lo que se complementa con el coeficiente $\symref{sym-cramer}{V}$ de Cramér como medida del tamaño del efecto.

\subsection{Visualización geoespacial}

Los mapas de calor, los mapas de puntos con grupos y los tableros permiten comunicar la distribución espacial y temporal de los incidentes a quienes toman decisiones operativas \citep{ristea2020crimegis,medinaariza2023hotspots}. En este estudio se emplearán para presentar los \textit{hotspots} identificados y para facilitar su interpretación por el personal de la municipalidad.

\section{Definición de términos básicos}

\begin{description}[leftmargin=0pt,labelindent=0pt,font=\bfseries]
\item[\termref{term-incident}{Ocurrencia de seguridad ciudadana}:] hecho atendido y registrado por el patrullaje municipal e integrado (serenazgo y Patrullaje Local Integrado), con fecha, turno, tipo y ubicación; no todos son delitos.
\item[\termref{term-spatial-cluster}{Grupo espacial (clúster)}:] conjunto de ocurrencias agrupadas por un algoritmo según su densidad.
\item[\termref{term-hotspot}{Hotspot}:] grupo espacial cuya concentración de ocurrencias supera la esperada por azar según un contraste estadístico.
\item[\termref{term-georeferencing}{Georreferenciación}:] asignación de coordenadas a un registro para ubicarlo en un sistema de referencia espacial.
\item[\termref{term-noise}{Ruido}:] punto que DBSCAN no asigna a ningún grupo por hallarse en una zona de baja densidad.
\item[\termref{term-shift}{Turno}:] franja horaria de operación del serenazgo (mañana, tarde y noche).
\item[Aleatoriedad espacial completa (\gls{csr}):] modelo en el que los puntos se distribuyen de forma independiente y uniforme en el área de estudio; sirve de referencia nula mínima; no incorpora población, vías ni patrullaje.
\end{description}
